\section{Related work}
\label{sec:related}

\paragraph{The residual stream and superposition.}
\citet{elhage2021mathematical} formalised the framing every cross-layer method
inherits: attention and MLP blocks read from and write to a shared additive
channel, so features persist through depth because the channel is a running
sum. \citet{elhage2022superposition} explains why single neurons are not the
unit of analysis: models encode more features than dimensions, in overlapping
directions. Together these motivate both the object of this study, the
per-layer write $\Dl$, and its principal statistical hazard --- consecutive
updates mechanically share a $\mp\vl$ term (\cref{sec:phase1:tautology}),
which we spend considerable effort separating from genuine structure.

\paragraph{Per-layer sparse dictionaries.}
SAEs trained on one layer's activations recover monosemantic features
\citep{cunningham2023sparse,bricken2023monosemanticity} and scale to frontier
models \citep{templeton2024scaling}. TopK SAEs replace the L1 penalty with a
hard sparsity constraint \citep{makhzani2013ksparse,gao2024scaling}. We use
public per-layer suites as frozen sources --- the MLSAE release for
Pythia \citep{lawson2024mlsae}, the reformatted GPT-2-small residual suite
\citep{bloom2024gpt2saes}, an open Qwen2 suite \citep{chanin2024qwensaes}, and
EleutherAI's Pythia-70m suite \citep{eleutherai2024sae} --- and additionally
train matched L1 and TopK dictionaries on our own cached activations, which
turns out to matter: the sparsity mechanism, not the model, controls whether
single-feature decoder directions are identifiable across seeds
(\cref{sec:evaluation:stability}).

\paragraph{Cross-layer dictionaries.}
Crosscoders fit one dictionary with per-layer decoders across all layers at
once \citep{lindsey2024crosscoders}. MLSAEs share the dictionary and let
activations vary per layer; their headline observation --- a latent is often
active at exactly one layer per token, but at \emph{different} layers across
tokens --- already strains the picture of a feature with a fixed home depth
\citep{lawson2024mlsae}. RouteSAE adds a router that picks which layer a
shared SAE decodes per token \citep{shi2025route}: a \emph{where-to-read}
mechanism, complementary to the \emph{how-updates-propagate} question here.
Jacobian SAEs sparsify the causal sensitivities across a single MLP block
\citep{farnik2025jacobian} --- causal rather than correlational cross-layer
structure, but with no assumption (or test) of depth-shared form.

\paragraph{Cross-layer transcoders.}
CLTs \citep{ameisen2025circuit} are the closest prior method and this paper's
primary conceptual baseline. A CLT feature reads at one source layer $\ell'$
and writes into every layer $\ell \ge \ell'$ through a separate, fully free
decoder per (source, target) pair --- $O(L^2)$ decoder matrices with no
parameter sharing along depth. Our H2 is precisely the hypothesis that those
matrices compress to a kernel in the depth difference $\ell-\ell'$; H2's
failure (\cref{sec:phase2}) is therefore an independent \emph{validation} of
CLT's unconstrained parameterisation, and our SSM-vs-CLT comparison
(\cref{sec:evaluation:clt}) asks the surviving quantitative question: how much
reconstruction do the free per-pair decoders buy over a first-order
depth-indexed state at matched sources?

\paragraph{Depth-wise lenses.}
The Tuned Lens learns per-layer affine probes that decode intermediate
residual states into the final-output distribution \citep{belrose2023tuned}.
We use it twice: its implied per-layer update
$\hat r_{\ell+1}-\hat r_\ell$ is the calibrated baseline our cheap H1
diagnostic must beat before the expensive tests run (it does,
\cref{sec:phase1:ridge}), and it reappears as an update-prediction baseline in
the evaluation, where --- like every method not trained on $\Dl$ --- it scores
below zero (\cref{sec:evaluation:baselines}).

\paragraph{Factual recall.}
We test on subject--relation--object recall because it is the
best-instrumented circuit family in the literature: causal localisation to
mid-depth MLPs \citep{meng2022rome,dai2022knowledge}, a staged
enrich-then-extract account \citep{geva2023dissecting}, and a standing caveat
that causal-tracing locations and editing locations dissociate
\citep{hase2023does}. Our data derive from CounterFact \citep{meng2022rome},
with Pararel \citep{elazar2021pararel} recorded in the protocol as the scaling
path we ultimately did not need (\cref{app:prereg}). Our distributed-code null
(\cref{sec:evaluation:causal}) is consistent with this literature's drift away
from single-component stories of recall.

\paragraph{Signal-processing motivation.}
Convolution Kernel Compensation decomposes high-density surface EMG under a
convolutive mixing model \citep{holobar2007multichannel}. The analogy ---
sources with characteristic propagation patterns along a depth-like axis ---
supplied H2's form (a shared kernel in the depth difference). It supplied
nothing else: every model fit here is ridge regression, LASSO, alternating
minimisation, or gradient descent on a linear SSM, and H2 was treated as a
hypothesis to kill, which the data did.

\paragraph{Infrastructure.}
Activations are collected with TransformerLens \citep{nanda2022transformerlens}
from Pythia \citep{biderman2023pythia}, GPT-2, and Qwen2 checkpoints; generic
text is drawn from the Pile \citep{gao2020pile}. Optimisation uses FISTA
\citep{beck2009fista} for the sparse-code updates; the Granger-style tests are
calibrated against a Freedman--Lane permutation scheme
\citep{freedman1983nonstochastic} in \cref{app:stats}.
